\chapter{总结与展望}
\label{chap:conclusion}

\section{全文总结}
\label{sec:conclusion_summary}

本文围绕复杂水下环境中的AFDM无线光通信关键技术展开研究，针对低动态多径场景与高动态双选择性场景分别提出了两类传输方案，并从系统建模、参数设计、检测方法和误码性能等方面进行了分析与验证。

针对低动态水下场景中散射主导的频率选择性衰落问题，本文提出了ACO-AFDM传输方案。该方案将ACO机制与AFDM调制相结合，在满足IM-DD架构非负实值约束的前提下，利用DAFT域啁啾结构增强系统对多径传播的适应能力。通过数学推导证明了限幅噪声仅分布于偶数子载波位置，保证了ACO机制在AFDM框架下的有效性；推导了啁啾参数 $c_1$ 的设计准则，建立了奇异点规避与多径分离条件；从等效信道矩阵结构角度揭示了ACO-AFDM非对角信道相较于 ACO-OFDM 对角信道的本质优势。仿真结果表明，ACO-AFDM 在双径信道下相较 ACO-OFDM 可获得约 2 dB 的性能增益，且在路径数从 2 增至 10 的鲁棒性测试中性能退化速度显著慢于传统光 OFDM 方案。

针对高动态水下场景中多普勒扩展与相位噪声的联合失真问题，本文提出了相干AFDM-IM传输方案。该方案在相干接收架构下引入AFDM波形与索引调制机制，建立了包含多径、多普勒频移和维纳相位噪声的统一系统模型。利用AFDM-IM静默位置的零结构先验，提出了基于泄露能量最小化的盲残余 CFO 补偿方法，在无额外导频开销下实现主导频偏的精准补偿，并对维纳相位噪声保持鲁棒容忍。通过代价函数全局收敛性分析证明了盲搜准则的可靠性，并定量推导了传统 20\% 导频开销所引起的约 0.97 dB 等效 SNR 惩罚。仿真结果表明，所提方案在高动态条件下显著优于基准 OFDM-IM，在归一化多普勒频移 0 至 0.20 的宽泛范围内均保持可靠通信能力。表~\ref{tab:thesis_contribution_summary} 对全文两类核心场景、对应架构选择、关键方法与主要性能结论进行了集中归纳，可作为本文技术主线的最终摘要。

在系统设计方面，本文建立了包含多径传播、多普勒频移以及维纳相位噪声的统一离散系统模型，对AFDM-IM信号在双选择性信道中的传输特性进行了分析。在此基础上，充分利用索引调制所产生的静默位置零结构先验信息，提出了一种基于泄露能量最小化准则的盲残余载波频偏补偿方法。该方法无需额外导频辅助，即可通过搜索最优频偏补偿参数实现主导频偏估计与校正，从而有效降低频偏失配带来的性能损失。

为了验证该方法的理论合理性，本文进一步对所构建的代价函数进行了分析，证明了其在真实频偏附近具有明确的全局最优特性，从理论上保证了盲搜索补偿方法的有效性。同时，针对传统导频辅助方案存在频谱资源损耗的问题，定量推导了20%导频开销对应的有效信噪比损失约为0.97 dB，说明盲补偿机制在频谱效率方面具有明显优势。

仿真结果表明，在归一化多普勒频移从0增长至0.20的宽范围条件下，所提出的AFDM-IM系统始终保持较好的误码性能；相比传统OFDM-IM方案，其抗多普勒能力和抗频偏能力均得到明显提升。同时，所提出的盲补偿方法能够有效抑制残余频偏引起的能量泄露现象，在不增加额外导频开销的情况下实现可靠检测，验证了方案在高动态复杂水下环境中的应用潜力。

总体而言，本文针对复杂水下环境中的两类典型通信场景，分别构建了面向低动态多径信道的ACO-AFDM方案以及面向高动态双选择性信道的AFDM-IM方案。两项研究均以AFDM波形为核心，通过结合不同接收架构和系统需求，实现了对多径衰落、多普勒扩展以及相位噪声等关键问题的针对性解决。研究结果表明，AFDM不仅能够有效提升系统的环境适应能力，而且具有良好的可扩展性和工程应用潜力。本文的工作为AFDM技术在水下无线光通信领域的进一步应用提供了理论基础和设计参考。

\section{未来展望}
\label{sec:future_work}

虽然本文针对复杂水下环境中的AFDM无线光通信技术开展了较为系统的研究，并取得了一定成果，但相关研究仍主要基于理论分析与计算机仿真验证，距离实际工程应用仍存在进一步深入研究的空间。未来可从以下几个方面继续开展研究工作。

在实验平台验证方面，后续可依托水下无线光通信实验系统开展端到端实物测试研究。通过在水槽环境、循环水环境以及开放水域环境中搭建实验链路，对本文提出的ACO-AFDM和AFDM-IM方案进行实际验证，进一步分析系统在真实水体中的性能表现。实验结果不仅能够验证理论模型与仿真结论的准确性，还能够为后续系统优化提供更加可靠的工程依据。

在信道建模方面，在信道建模方面，未来可进一步建立更加符合实际应用场景的综合信道模型。本文主要考虑了散射、多径传播、多普勒扩展以及相位噪声等因素，而实际系统中还可能受到动态指向误差、平台振动、光束漂移、海洋湍流以及收发端失准等因素的影响。将这些因素统一纳入建模框架，有助于更加准确地描述真实水下光通信信道特性，并提高算法设计的实际适用性。

在低复杂度实现方面，可围绕近似检测、迭代检测以及基于学习的检测方法开展研究，在保持性能的前提下降低接收端计算开销，使算法更接近实际接收机的处理能力。

在系统拓展方面，可研究AFDM与多输入多输出、中继协作及自适应资源分配机制的结合，将方案从点对点单链路拓展至多节点网络化传输框架；也可探索在不同链路状态下，根据环境条件在不同波形参数与接收策略之间进行自适应切换的统一设计方法。
最后，随着智能海洋和自主水下系统的发展，未来还可研究AFDM技术与人工智能算法的深度融合。例如利用机器学习方法实现信道预测、参数优化和智能资源调度，构建具有环境感知与自主决策能力的新型水下通信系统，从而进一步提升复杂海洋环境下的通信效率和系统智能化水平。

综上所述，AFDM在水下无线光通信领域仍具有广阔的发展空间。随着信道建模、实验验证、智能检测以及网络化应用等方向的不断推进，其有望成为未来高可靠、高速率水下光通信系统的重要候选技术之一。